\section{Aksesibilitas dalam CSS}

CSS memainkan peran penting dalam aksesibilitas web---bukan hanya mempercantik tampilan, tetapi juga memastikan konten dapat digunakan oleh semua orang, termasuk pengguna dengan disabilitas visual, motorik, atau kognitif. Beberapa praktik CSS langsung mempengaruhi pengalaman pengguna assistive technology \cite{mdnAccessibility}.

\begin{table}[htbp]
\centering
\small
\begin{tabular}{|P{3.2cm}|P{2.4cm}|P{3cm}|}
\hline
\textbf{Aspek} & \textbf{Teknik CSS} & \textbf{Manfaat} \\
\hline
Indikator fokus & \code{:focus}, \code{:focus-visible}, outline/box-shadow & Pengguna keyboard tahu posisi fokus \\
\hline
Reduced motion & \parbox[t]{2.4cm}{\raggedright\code{@media (prefers-}\linebreak[0]\code{reduced-motion)}} & Kurangi animasi bagi yang sensitif \\
\hline
Kontras warna & Ratio $\geq$ 4.5:1 teks, $\geq$ 3:1 besar & Teks terbaca di semua kondisi \\
\hline
Skip to content & \code{.sr-only} class & Link cepat ke konten utama \\
\hline
Responsive text & Hindari \code{px} untuk font-size & Teks mengikuti preferensi pengguna \\
\hline
\end{tabular}
\caption{Teknik aksesibilitas CSS dan manfaatnya}
\end{table}

\begin{webcode}[language=CSS, caption={CSS aksesibilitas: focus, reduced motion, sr-only}]
/* Focus visible */
:focus-visible {
  outline: 2px solid #2a6fb0;
  outline-offset: 2px;
}
/* Reduced motion */
@media (prefers-reduced-motion: reduce) {
  *, *::before, *::after {
    animation-duration: 0.01ms !important;
    transition-duration: 0.01ms !important;
  }
}
/* Visually hidden (screen reader only) */
.sr-only {
  position: absolute; width: 1px; height: 1px;
  padding: 0; margin: -1px; overflow: hidden;
  clip: rect(0,0,0,0); border: 0;
}
\end{webcode}

\begin{konsep}
\textbf{Jangan Sembunyikan Fokus.} Aturan \code{*:focus \{ outline: none; \}} adalah salah satu kesalahan aksesibilitas paling umum. Gunakan \code{:focus-visible} untuk menampilkan outline hanya saat navigasi keyboard (bukan saat klik mouse). Jika outline default tidak sesuai desain, ganti dengan \code{box-shadow} atau \code{border}---jangan pernah hilangkan tanpa pengganti yang terlihat \cite{webdevAccessibility}.
\end{konsep}

